\documentclass[tikz,border=6pt]{standalone}
\usepackage{fontspec}
\usepackage[mathrm=sym]{unicode-math}
\setmainfont{texgyrepagella-regular.otf}[
    BoldFont = texgyrepagella-bold.otf,
    ItalicFont = texgyrepagella-italic.otf,
    BoldItalicFont = texgyrepagella-bolditalic.otf
]
\setmathfont{texgyrepagella-math.otf}[AutoFakeBold = 2.5]
\usepackage{tikz}
\usetikzlibrary{arrows.meta, decorations.markings}

\begin{document}
\begin{tikzpicture}[>=Stealth, line cap=round, line join=round,
  arrowed/.style={
    decoration={
      markings,
      mark=at position 0.55 with {\arrow{Stealth[length=5pt, width=3.5pt]}}
    },
    postaction={decorate}
  },
  arrowedpos/.style={
    decoration={
      markings,
      mark=at position #1 with {\arrow{Stealth[length=5pt, width=3.5pt]}}
    },
    postaction={decorate}
  }
]

  % ========================================================
  % ROW 1: Bethe-Salpeter equation in pp channel
  % ========================================================
  \begin{scope}[shift={(0, 1.4)}]

    % LHS: pp 4-point function (circle)
    \begin{scope}[shift={(-4.8, 0)}]
      \draw[thick, line width=1pt, arrowed] (-1.1, 0.35) -- (-0.42, 0.35);
      \draw[thick, line width=1pt, arrowed] (-1.1, -0.35) -- (-0.42, -0.35);
      \draw[thick, line width=1pt, arrowed] (0.42, 0.35) -- (1.1, 0.35);
      \draw[thick, line width=1pt, arrowed] (0.42, -0.35) -- (1.1, -0.35);
      \draw[thick, fill=white] (0,0) circle (0.46);
    \end{scope}

    % Equals sign
    \node[font=\large] at (-3.2, 0) {$=$};

    % Term 1: Direct free pp propagators
    \begin{scope}[shift={(-2.0, 0)}]
      \draw[thick, line width=1pt, arrowed] (-0.7, 0.35) -- (0.7, 0.35);
      \draw[thick, line width=1pt, arrowed] (-0.7, -0.35) -- (0.7, -0.35);
    \end{scope}

    % Plus sign
    \node[font=\large] at (-0.8, 0) {$+$};

    % Term 2: Exchange free pp propagators
    \begin{scope}[shift={(0.4, 0)}]
      \draw[thick, line width=1pt, arrowedpos=0.4] (-0.7, 0.35) -- (0.7, -0.35);
      \draw[thick, line width=1pt, arrowedpos=0.4] (-0.7, -0.35) -- (0.7, 0.35);
    \end{scope}

    % Plus sign
    \node[font=\large] at (1.6, 0) {$+$};

    % Term 3: K^{pp} + intermediate lines + circle
    \begin{scope}[shift={(2.5, 0)}]
      % Incoming legs to K^{pp}
      \draw[thick, line width=1pt, arrowed] (0, 0.35) -- (0.7, 0.35);
      \draw[thick, line width=1pt, arrowed] (0, -0.35) -- (0.7, -0.35);

      % Box K^{pp}
      \draw[thick, fill=white] (0.7, -0.42) rectangle (1.7, 0.42);
      \node[font=\normalsize] at (1.2, 0) {$K^{\mathrm{pp}}$};

      % Intermediate legs
      \draw[thick, line width=1pt, arrowed] (1.7, 0.35) -- (2.58, 0.35);
      \draw[thick, line width=1pt, arrowed] (1.7, -0.35) -- (2.58, -0.35);

      % Circle
      \draw[thick, fill=white] (3.0, 0) circle (0.46);

      % Outgoing legs
      \draw[thick, line width=1pt, arrowed] (3.42, 0.35) -- (4.2, 0.35);
      \draw[thick, line width=1pt, arrowed] (3.42, -0.35) -- (4.2, -0.35);
    \end{scope}

  \end{scope}

  % ========================================================
  % ROW 2: Expansion of irreducible pp kernel K^{pp}
  % ========================================================
  \begin{scope}[shift={(0, -1.2)}]

    % LHS: Box K^{pp}
    \begin{scope}[shift={(-4.8, 0)}]
      \draw[thick, line width=1pt, arrowed] (-1.1, 0.35) -- (-0.5, 0.35);
      \draw[thick, line width=1pt, arrowed] (-1.1, -0.35) -- (-0.5, -0.35);
      \draw[thick, line width=1pt, arrowed] (0.5, 0.35) -- (1.1, 0.35);
      \draw[thick, line width=1pt, arrowed] (0.5, -0.35) -- (1.1, -0.35);

      \draw[thick, fill=white] (-0.5, -0.42) rectangle (0.5, 0.42);
      \node[font=\normalsize] at (0, 0) {$K^{\mathrm{pp}}$};
    \end{scope}

    % Approx sign
    \node[font=\large] at (-3.2, 0) {$\simeq$};

    % Term 1: 4-point vertex
    \begin{scope}[shift={(-2.0, 0)}]
      \coordinate (C) at (0, 0);
      \draw[thick, line width=1pt, arrowedpos=0.6] (-0.7, 0.4) to[out=0, in=140] (C);
      \draw[thick, line width=1pt, arrowedpos=0.6] (-0.7, -0.4) to[out=0, in=-140] (C);
      \draw[thick, line width=1pt, arrowedpos=0.6] (C) to[out=40, in=180] (0.7, 0.4);
      \draw[thick, line width=1pt, arrowedpos=0.6] (C) to[out=-40, in=180] (0.7, -0.4);
      \filldraw[black] (C) circle (2pt);
    \end{scope}

    % Plus sign
    \node[font=\large] at (-0.8, 0) {$+$};

    % Term 2: particle-hole loop between pp lines
    \begin{scope}[shift={(0.8, 0)}]
      \coordinate (VT) at (-0.35, 0.35);
      \coordinate (VB) at (0.35, -0.35);

      % Top straight line
      \draw[thick, line width=1pt, arrowedpos=0.35] (-1.0, 0.35) -- (VT);
      \draw[thick, line width=1pt, arrowedpos=0.65] (VT) -- (1.0, 0.35);

      % Bottom straight line
      \draw[thick, line width=1pt, arrowedpos=0.35] (-1.0, -0.35) -- (VB);
      \draw[thick, line width=1pt, arrowedpos=0.65] (VB) -- (1.0, -0.35);

      % Loop between VT and VB: circulating loop
      % Left arc: pointing up-right from VB to VT
      \draw[thick, line width=1pt, arrowed] (VB) to[out=160, in=-100] (VT);
      % Right arc: pointing down-right from VT to VB
      \draw[thick, line width=1pt, arrowed] (VT) to[out=-20, in=80] (VB);

      \filldraw[black] (VT) circle (2pt);
      \filldraw[black] (VB) circle (2pt);
    \end{scope}

  \end{scope}

\end{tikzpicture}
\end{document}
